% ── Section V: Results & Discussion ──────────────────────────────────
% Conference budget: ~1.5 pages. WARNING: this section roughly doubled
% when the real numbers landed and now holds four tables. Measure the
% actual page count before submission; if it must be cut, the cheapest
% losses are (a) tab:detection -> two sentences of inline numbers,
% (b) tab:bias -> keep only K=0 and K=29, (c) the middle three rows of
% tab:ablation_results -> a single sentence, though that costs the
% gate/NSA interaction argument, which is the most novel claim here.
% 5.1 is REAL already (synthetic studies, reproducible from the repo).
% 5.2–5.5 await VN footage — structure and analysis prompts in place.
% ─────────────────────────────────────────────────────────────────────

\section{Results and Discussion}
\label{sec:results}

\subsection{Controlled Synthetic Studies: Motion Model}
\label{subsec:synthetic}

Before end-to-end evaluation, we isolate the motion model on synthetic
trajectories with known ground truth (generator included in our code).
Three findings:
(i)~on constant-turn trajectories, CTRV reduces one-step prediction
RMSE by \textbf{48\%} relative to the linear Kalman baseline, and by
\textbf{13\%} on sinusoidal weaving, with no regression on straight
lines;
(ii)~across a 12-frame occlusion in the middle of a turn, the CTRV
track coasts along the arc and re-associates with its original
identity, while the linear baseline drifts off the arc and
re-identifies (1 IDSW vs.\ 0);
(iii)~the failure boundary is characterized: when an occlusion begins
immediately after a turn-rate sign flip, CTRV extrapolates the stale
$\omega$ and a straight coast is closer to the truth --- motivating the
context-gated activation of CTRV (C1) rather than unconditional
replacement.

\subsection{Ablation Study}
\label{subsec:ablation_results}

Table~\ref{tab:ablation_results} reports the ablation on the
human-annotated \textsc{dense} segment, all rows consuming identical
detections so that every difference is attributable to tracking.
C2 is the decisive contribution: on motorcycles it raises IDF1 from
$0.844$ to $0.933$ and removes $85\%$ of identity switches
($189 \rightarrow 29$).
C1 does not improve on a well-chosen fixed threshold here, and C3 leaves
the tracking metrics unchanged by construction; we examine both below
rather than presenting them as gains.

\begin{table}[h]
  \centering
  \caption{Ablation on the annotated \textsc{dense} segment (600 frames,
           $\sim$47 tracks/frame), motorcycle class. Identical
           detections for every row. \textbf{Bold} = best.}
  \label{tab:ablation_results}
  \begin{tabular}{lcccc}
    \toprule
    Config & MOTA$\uparrow$ & IDF1$\uparrow$ & IDSW$\downarrow$
           & Frag$\downarrow$ \\
    \midrule
    Exp~0 (baseline)          & 0.843 & 0.844 & 189 & 199 \\
    \quad + CTRV              & 0.857 & 0.850 & 158 & 175 \\
    \quad + CTRV, NSA         & 0.869 & 0.869 & 123 & 131 \\
    \quad + CTRV, gate        & 0.854 & 0.868 & 182 & 157 \\
    Exp~1 (+C2, full)         & \textbf{0.893} & \textbf{0.933}
                              & \textbf{29} & \textbf{64} \\
    Exp~2 (+C1, C2)           & 0.878 & 0.909 & 99 & 97 \\
    Exp~3 (+C1, C2, C3)       & 0.878 & 0.909 & 99 & 97 \\
    \bottomrule
  \end{tabular}
\end{table}

\paragraph{The parts of C2 are not additive}
The three middle rows isolate the interaction argued in
Sec.~\ref{subsec:c2}.
Uncertainty gating \emph{alone} raises identity switches above plain
CTRV ($182$ vs.\ $158$): without confidence-scaled measurement noise the
innovation ellipse is wide and mis-centered, so rescued matches
frequently attach a lost track to a neighboring motorcycle.
NSA alone is a modest gain ($158 \rightarrow 123$).
Together they reach $29$ --- far beyond the sum of the parts.
This is why we present C2 as one mechanism: an implementation that
adopted the gate without the noise model would measure worse than the
baseline it was meant to improve.

\paragraph{Detection quality}
Table~\ref{tab:detection} isolates what tracking contributes over
per-frame detection.
At the operating threshold the detector alone reaches
$P{=}0.890$, $R{=}0.847$; after tracking, $P{=}0.950$ and $R{=}0.968$.
Both rise together, which no confidence threshold can achieve --- moving
the cut-off only slides along the precision--recall curve
($0.30 \rightarrow 0.70$ buys $+0.20$ precision for $-0.33$ recall).
Temporal association recovers objects the detector drops in individual
frames while discarding detections that fail to persist, leaving the
per-frame curve rather than moving along it.

\begin{table}[h]
  \centering
  \caption{Detection quality on the \textsc{dense} segment.
           mAP@0.5 over all detections gives the detector ceiling.}
  \label{tab:detection}
  \begin{tabular}{lcccc}
    \toprule
    Source & P$\uparrow$ & R$\uparrow$ & F1$\uparrow$ & mAP@0.5 \\
    \midrule
    Detector, all detections & 0.497 & 0.984 & 0.661 & 0.925 \\
    Detector @ $c \geq 0.30$ & 0.746 & 0.943 & 0.833 & -- \\
    Detector @ $c \geq 0.50$ & 0.890 & 0.847 & 0.868 & -- \\
    Detector @ $c \geq 0.70$ & 0.946 & 0.615 & 0.745 & -- \\
    \midrule
    Baseline tracker         & 0.941 & 0.944 & 0.942 & -- \\
    Proposed (C2)            & \textbf{0.950} & \textbf{0.968}
                             & \textbf{0.959} & -- \\
    \bottomrule
  \end{tabular}
\end{table}

\paragraph{Ground-truth over-merging (open issue)}
An automatic consistency pass finds that $74$ of $335$ annotated
identities ($22\%$) contain a join across which the box teleports or
rescales implausibly --- several vehicles chained into one identity.
Two were verified manually.
This error is not neutral: an over-merged identity rewards a tracker that
bridges the gap between two vehicles and penalises one that correctly
refuses, which is precisely the behaviour C2's gate supplies.
Splitting the flagged joins and re-scoring moves the proposed margin from
$+0.089$ to between $+0.061$ and $+0.014$ motorcycle IDF1, with the
baseline \emph{improving} and the proposed method \emph{degrading}; plain
CTRV improves while CTRV+gate degrades, isolating the gate as the
component the annotation error rewards.
The proposed configuration leads on every version of the ground truth we
construct, but we regard the magnitude in
Table~\ref{tab:ablation_results} as an upper bound pending a re-annotation
pass targeting these joins, and we report it as such rather than as a
settled figure.

\paragraph{Bounding the pre-label bias}
Because the ground truth was seeded from the proposed tracker's output
(Sec.~\ref{subsec:datasets}), fragmentations the annotator missed persist
as agreement and flatter that row.
We bound this rather than merely disclosing it.
An automatic consistency pass over the finished ground truth flags $29$
residual merge candidates; applying the top $K$ of them and re-scoring
every configuration gives Table~\ref{tab:bias}.
Correcting \emph{all} $29$ --- a deliberate worst case, since inspection
shows only the highest-confidence third to be genuine --- moves the
proposed margin over baseline from $+0.089$ to $+0.085$ IDF1.
Both rows fall together, because a merged identity penalizes whichever
configuration split it, and the baseline splits in more places.

\begin{table}[h]
  \centering
  \caption{Sensitivity of the proposed--baseline margin to correcting
           $K$ residual ground-truth fragmentations (motorcycle IDF1).}
  \label{tab:bias}
  \begin{tabular}{lccc}
    \toprule
    $K$ corrected & GT identities & Baseline & Proposed \\
    \midrule
    0 (as annotated) & 335 & 0.844 & 0.933 \\
    3                & 332 & 0.844 & 0.931 \\
    10               & 325 & 0.842 & 0.928 \\
    29 (all flagged) & 306 & 0.831 & 0.916 \\
    \bottomrule
  \end{tabular}
\end{table}

\subsection{Where C1 and C3 Stand}
\label{subsec:c3_limits}

We report two negative results, both of which we consider more useful
than the numbers they replace.

\paragraph{C1 is untested, not refuted}
After calibration (Table~\ref{tab:param_table}) C1 reaches MOTA $0.905$
against $0.917$ for the fixed configuration --- it does not help.
But the calibrated optimum coincides with the fixed operating point,
which is the signature of a segment that never changes regime: density
on \textsc{dense} never falls below the switching threshold, so C1 holds
one condition for $96.5\%$ of frames.
An adaptation mechanism cannot be evaluated where there is nothing to
adapt to.
The honest statement is that C1's shipped tuning was harmful, that
calibration removes the harm, and that the premise remains untested
pending a segment that straddles the switching boundary.

\paragraph{C3's recall branch cannot help, for a measurable reason}
Raising $\tau_{\mathrm{mappo}}$ monotonically degrades MOTA
($0.905 \rightarrow 0.705$ at $\tau = 0.80$).
This is not a tuning failure but a consequence of two measurements:
the reliability score separates true from false tracks only weakly
(ROC AUC $0.609$), while false positives are just $5.5\%$ of tracks
after C2.
At its most enriched cut-off the filter still discards four true tracks
for every false one, so no threshold trades favorably.
Notably C2 causes this: by removing most spurious tracks it leaves C3
little to remove, so the contributions interact rather than compose.
Improving C3 means improving the score --- AUC $0.609$ is the quantity
to move --- not the threshold.
The asymmetry C3 was designed around may still hold on the accident
branch, where discarding four true tracks to suppress one false alert is
a defensible trade; that requires accident footage we do not yet have.

\subsection{Comparison with Prior Systems}
\label{subsec:comparison}

Table~\ref{tab:comparison} contrasts capabilities and, where
measurable, tracking accuracy.
Hikvision AID's internals are unpublished, so it appears in the
capability matrix only.
The YOLO-only row quantifies temporal confirmation: without it, no
confidence threshold reaches an F1 above $0.868$
(Table~\ref{tab:detection}), against $0.959$ once tracks are formed.
The ByteTrack row isolates our contribution proper, since it shares the
detector and the association logic and differs only in the motion
mechanism of Sec.~\ref{subsec:c2}.

\begin{table}[h]
  \centering
  \caption{Capability and accuracy comparison.}
  \label{tab:comparison}
  \resizebox{\linewidth}{!}{%
  \begin{tabular}{lccccc}
    \toprule
    Method & Mixed & Class-cond. & Adaptive & Severity & MOTA$\uparrow$ \\
    \midrule
    Hikvision AID~\cite{Hikvision_AID}   & --         & --         & --         & --         & n/a \\
    YOLO-only                            & \checkmark & --         & --         & --         & n/a \\
    ByteTrack~\cite{Zhang2022_bytetrack} & \checkmark & --         & --         & --         & 0.878 \\
    Proposed                             & \checkmark & \checkmark & \checkmark & \checkmark & \textbf{0.917} \\
    \bottomrule
  \end{tabular}}
\end{table}

\subsection{Qualitative Results}
\label{subsec:qualitative}

\todo{Two figures once VN footage is processed:
(a) side-by-side weaving sequence, baseline IDs fragmenting vs.\
proposed IDs stable; (b) motivation frames --- baseline false alert and
lost motorcycle track, annotated.}

\subsection{Runtime}
\label{subsec:runtime}

Table~\ref{tab:runtime} reports tracker-side latency with the detector
excluded, isolating what the proposed modules cost.
C2 adds $0.85$\,ms per frame ($+22\%$) over the linear-filter baseline,
and the ordering follows the computation: the EKF is the bulk of it,
NSA is a single scaling of $\mathbf{R}$ and costs about two points, and
the gate adds one Mahalanobis test per lost track.
% REWRITTEN 27/8: the previous version asserted "a YOLO11 inference cost of
% tens of milliseconds" and a 33 ms budget. Both were wrong -- the shipped
% detector costs 130.7 ms and the source is 25 fps (40 ms). The verdict on C2
% happened to survive, which is exactly why the premise had to be fixed rather
% than left standing. See experiments/realtime/evidence.md.
C2 is not the term that decides real-time feasibility, but the margin is
narrower than tracker-side numbers suggest.
Measured end to end on the same machine, the shipped detector
(a YOLO11n \emph{segmentation} checkpoint at $640$\,px) costs
$130.7$\,ms per frame --- $91\%$ of a $143.2$\,ms frame, against
$10.1$\,ms for tracking (Table~\ref{tab:endtoend}).
At $7.0$\,fps the pipeline misses the $40$\,ms budget set by the $25$\,fps
source by $3.6\times$, so the operating point, not the tracker, is what
stands between this system and real time.
Section~\ref{sec:implementation} takes that apart: the segmentation head
alone accounts for a $2.25\times$ factor over a detection head of the same
scale for masks the pipeline never reads, and a detection head at $448$\,px
reaches $37.6$\,fps within budget.

\begin{table}[h]
  \centering
  \caption{Tracker-side per-frame latency, detector excluded
           (13\,726 frames, CPU). ``p95 dense'' is the worst-case
           window at $\sim$47 tracks/frame.}
  \label{tab:runtime}
  \begin{tabular}{lcccc}
    \toprule
    Config & mean (ms) & p50 & p95 & p95 dense \\
    \midrule
    Baseline (linear KF)     & 3.93 & 3.57 & 6.68 & 9.67 \\
    \quad + CTRV             & 4.44 & 3.96 & 7.51 & 9.87 \\
    \quad + CTRV, NSA        & 4.51 & 4.03 & 7.68 & 9.82 \\
    \quad + CTRV, gate       & 4.67 & 4.04 & 8.15 & 13.68 \\
    C2 (full)                & 4.78 & 4.16 & 8.28 & 14.09 \\
    C1 + C2 + C3             & 5.42 & 4.67 & 9.71 & 13.80 \\
    \bottomrule
  \end{tabular}
\end{table}

Each configuration was timed in a separate process.
Timing all of them in one run makes the measurement track position
rather than variant: the same baseline code measured $13.19$\,ms first
and $4.07$\,ms last, a $3.2\times$ spread from warm-up and scheduling
that exceeds the differences being measured and flatters whichever
configuration runs last.
